% Add \usepackage[labelsep=space]{caption} to the document preamble

\section{Preliminaries}
\label{secpreliminaries}

A discovery task $d$ defines the scientific objects that may be proposed, the evidence required to evaluate them, and the resources available for search. Each completed allocation runs one search branch and returns a candidate $x$ with an immutable evidence packet $\eta$. An independent verifier decides whether the candidate is admissible and assigns a bounded verified score $R_d(x,\eta)\in[L,U]$. Scores are oriented so that larger values are better. When evaluation is noisy, $R_d$ is a prespecified lower confidence value rather than a single measurement.

Let $C_i$ be the cost of the $i$th completed branch. Under total budget $B$, the final completed branch is
\begin{equation}
\tau_B
=
\max\left\{
t
\ \middle|\
\sum_{i=1}^{t}C_i\leq B
\right\}.
\label{eqstopping}
\end{equation}
The verified record after $t$ branches is
\begin{equation}
M_t
=
\max\left(
M_0,
\max_{\substack{1\leq i\leq t\\x_i\text{ is admitted}}}
R_d(x_i,\eta_i)
\right).
\label{eqrecord}
\end{equation}
EVOLVE seeks a policy $\Pi$ that maximizes the expected final record
\begin{equation}
J_B(\Pi)
=
\mathbb E_{\Pi}\left[M_{\tau_B}\right].
\label{eqobjective}
\end{equation}
Admissibility, scientific diversity, and resource limits guide which branches may be run. They are not added to the verified score. This keeps the search focused on the best confirmed scientific result.

\section{Framework}
\label{secframework}

We call the framework EVOLVE for Evidence-Verified Option Learning and Value Exploration. EVOLVE organizes discovery as repeated branching from a verified scientific archive. A scheduler chooses where to begin, which agent role to use, which multi-step search tactic to run, which harness version should support it, and how much budget it receives. An independent verifier then returns evidence that updates the archive, the scheduler, the option memory, and the selected agent policy.

The default system uses one frozen language model backbone with three separate LoRA adapters. The adapters implement the Scout, Mechanist, and Challenger roles. Each role has its own optimizer state, random stream, and private working memory. This gives EVOLVE three evolving policies while keeping the shared foundation model fixed.

\begin{figure}[t]
\centering
\resizebox{\linewidth}{!}{\begin{tikzpicture}[
  font=\sffamily\small,
  node distance=5mm and 10mm,
  box/.style={draw, rounded corners=2pt, very thick, align=center, minimum height=8mm, fill=white},
  flow/.style={-{Latex[length=2mm]}, thick},
  input/.style={box, draw=blue!65!black, fill=blue!5, minimum width=35mm},
  control/.style={box, draw=violet!65!black, fill=violet!6, minimum width=47mm},
  agent/.style={box, draw=teal!65!black, fill=teal!6, minimum width=66mm},
  process/.style={box, draw=orange!70!black, fill=orange!7, minimum width=47mm},
  verify/.style={box, draw=green!55!black, fill=green!7, minimum width=47mm},
  update/.style={box, draw=red!60!black, fill=red!5, minimum width=34mm}
]

\node[input] (archive) {Verified scientific archive};
\node[input, right=of archive] (harness) {Versioned harness pool};

\node[control, below=of archive, xshift=22.5mm] (scheduler) {Posterior record-improvement scheduler};
\node[agent, below=of scheduler] (agents) {One frozen backbone with three role-specific LoRAs\\Scout \quad Mechanist \quad Challenger};
\node[process, below=of agents] (branches) {Independent multi-step option branches};
\node[verify, below=of branches] (verifier) {Common independent verifier};

\node[update, below left=of verifier, xshift=-13mm] (archiveupdate) {Archive and\\record update};
\node[update, below=of verifier] (memory) {Causal option\\memory update};
\node[update, below right=of verifier, xshift=13mm] (learning) {Role LoRA\\learning update};

\draw[flow] (archive) -- (scheduler);
\draw[flow] (harness) -- (scheduler);
\draw[flow] (scheduler) -- (agents);
\draw[flow] (agents) -- (branches);
\draw[flow] (branches) -- (verifier);
\draw[flow] (verifier) -- (archiveupdate);
\draw[flow] (verifier) -- (memory);
\draw[flow] (verifier) -- (learning);

\node[font=\sffamily\footnotesize, below=3mm of memory] {Updates enter the next synchronized epoch};

\end{tikzpicture}
}
\caption{The main EVOLVE cycle. Archive states and harness versions enter the scheduler. A selected role runs an option branch under one frozen backbone. Verified evidence then updates the archive and record, causal option memory, and the corresponding role LoRA for the next epoch.}
\label{figframework}
\end{figure}

\subsection{EVOLVE at a Glance}
\label{secoverview}

Search is divided into epochs. An epoch is the interval between two synchronization barriers. The following sequence gives the complete operating loop.

\begin{enumerate}
\item Freeze the current role adapters, archive map, harness versions, verifier version, and record threshold.
\item Select promising and under-tested scientific states from the archive.
\item Form candidate allocations by combining a state, role, option, harness version, horizon, and cost.
\item Reserve small budget shares for randomized audits, archive coverage, all three roles, and new harness trials. Allocate the remaining budget by expected record improvement.
\item Run independent production branches and matched audit pairs.
\item Verify every completed branch and optionally send a suitable provisional candidate through bounded refinement.
\item At the synchronization barrier, update the record, archive, scheduler, causal option memory, and the LoRA adapter of each role from its own verified data.
\end{enumerate}

The same evidence is used in different ways. Production branches estimate where the next record is likely to come from. Randomized audit pairs estimate whether an option or harness caused an improvement relative to a matched alternative. Homogeneous groups of completed branches provide the rewards for test-time policy learning.

\subsection{Scientific States and Option Branches}
\label{secstatesoptions}

The archive groups candidates by scientific behavior rather than by wording or surface form. A versioned descriptor $\phi_v(x,\eta)$ summarizes measured responses, mechanistic structure, and active constraints. A mapping $q_v$ assigns the descriptor to a scientific state cell
\begin{equation}
c=q_v\bigl(\phi_v(x,\eta)\bigr).
\label{eqcell}
\end{equation}
Text similarity may help retrieval, while the scientific descriptor determines cell identity. The descriptor version remains fixed during an epoch so that evidence collected in the same epoch is comparable.

A cell keeps several useful candidates instead of retaining just one. These include its strongest confirmed candidate, a candidate with promising descendant value, and a small set of stepping stones. Local competition compares candidates within nearby scientific states. The best confirmed candidate across the whole archive is also retained as the global record holder. A coverage budget gives empty and under-tested cells regular opportunities without changing the scientific score.

This archive follows the local competition idea of quality-diversity search while allowing uncertain descriptors and more than one useful candidate per cell \citep{mouret2015,flageat2023}.

An option is a reusable multi-step search tactic
\begin{equation}
o=(I_o,\pi_o,\beta_o,H_o).
\label{eqoption}
\end{equation}
Here $I_o$ describes where the tactic may begin, $\pi_o$ is its internal proposal policy, $\beta_o$ is its stopping rule, and $H_o$ is its largest branch budget. Examples include exploring a new representation, modifying an active mechanism, challenging an assumption, or refining a provisional candidate.

This definition adapts the options formalism from temporal abstraction to scientific search \citep{sutton1999,bacon2017}.

Launching an option from a cell creates a local descendant tree. Different trees may later reach the same scientific cell, so the persistent search structure is an archive-linked provenance graph rather than one global tree. The tree records how a result was reached. Allocation decisions are made by the scheduler described below.

For a branch with horizon $H$, let $Z_H$ be the best verified score found by that branch or any of its descendants. The branch gain is
\begin{equation}
G_H
=
\left[Z_H-M_e^{\mathrm{start}}\right]_+,
\label{eqbranchgain}
\end{equation}
where $M_e^{\mathrm{start}}$ is the record frozen at the beginning of the epoch. This gives useful credit to a tactic that reaches a productive region even when its first proposal does not improve the record.

\subsection{Three Isolated Agent Roles}
\label{secagents}

EVOLVE uses one frozen backbone and three role-specific LoRA adapters.

\paragraph{Scout}
The Scout searches remote and underexplored cells. It proposes new option families, representations, and directions that may broaden the scientific search space.

\paragraph{Mechanist}
The Mechanist works from current constraints, mechanisms, and invariants. It develops targeted options whose value may appear over several descendant steps.

\paragraph{Challenger}
The Challenger searches for counterexamples, boundary cases, and minimal changes that test a current proposal. Its candidates are evaluated by the same independent verifier used for the other roles.

Each role keeps a private working memory and retrieval view. Raw transcripts remain private to the role that produced them. After verification, compact evidence packets become available to the other roles at the next synchronization barrier. Gradients update the corresponding role adapter and are not merged across roles. The result is a state-isolated ensemble rather than a claim that the three roles are statistically independent.

\subsection{Adaptive Allocation and Harness Selection}
\label{secscheduler}

A harness is the operating setup around an agent. It specifies the archive view, prompt and context assembly, role routing, branch breadth and depth, resource envelope, and verification schedule. EVOLVE does not rely on one harness being best in every context. Instead, the scheduler treats the harness as part of the allocation.

This choice is motivated by evidence that discovery harnesses can rank differently across models and tasks \citep{gupta2026}.

Let $h_v$ denote a harness and its version. A harness version remains fixed inside a branch. Any change to its context, routing, search schedule, memory access, or verification timing creates a new version that enters through a small budget-matched trial allocation. This makes the harness adaptive across branches and stable within each branch.

An allocation arm is
\begin{equation}
r=(c,a,o,h_v,H,C),
\label{eqarm}
\end{equation}
where $c$ is the starting cell, $a$ is the role, $o$ is the option, $h_v$ is the harness version, $H$ is the branch horizon, and $C$ is the total cost. Let $Z_r$ be the verified descendant maximum predicted for arm $r$. Given remaining budget $b_t$, the scheduler selects a portfolio $S_t$ by
\begin{equation}
S_t
\in
\arg\max_{S\,\mid\,\sum_{r\in S}C_r\leq b_t}
\mathbb E\left[
\left[
\max_{r\in S}Z_r-M_t
\right]_+
\ \middle|\
\mathcal D_t
\right].
\label{eqportfolio}
\end{equation}
This is posterior expected record improvement. It favors portfolios that are likely to beat the current record. The joint prediction also reduces the value of selecting several branches that are individually promising but likely to return the same result.

Before this performance-based allocation, fixed budget shares support randomized audits, empty or under-tested cells, all three roles, and provisional harness versions. The remaining budget follows Eq.~\eqref{eqportfolio}. Production outcomes update the tail posterior used by the scheduler. Admission probability and the score distribution of admitted candidates are modeled separately because both affect the chance of a new record.

\subsection{Verification, Feedback, and Causal Memory}
\label{secmemory}

All harnesses use the same external verifier. The verifier checks admissibility, evaluates the primary scientific measure, applies controlled perturbations when appropriate, and repeats measurements when outcomes are stochastic. It returns a compact certificate containing the admission decision, verified score, scientific descriptor, useful diagnostics, and branch lineage. Information reserved for blinded evaluation remains hidden from the agents.

EVOLVE uses three forms of memory. The archive stores verified scientific states. Each role keeps private working memory for its current reasoning. A causal option registry stores reusable evidence about options and harnesses. This registry contains structured records rather than unrestricted summaries. A record identifies the intervention, matched alternative, starting context, assignment probability, cost, horizon, verified outcome, uncertainty, and verifier version.

Production data alone may overrepresent options that the scheduler already prefers. A fixed share of the budget is therefore used for randomized audit pairs. Both branches start from the same state and frozen role checkpoint. They receive the same horizon and cost, while one branch receives the option and the other receives a matched continuation $o_0$. Their outputs remain separate until both branches close. The local branch effect is
\begin{equation}
\tau_H(c,a,h_v,o)
=
\mathbb E\left[
G_H(o)-G_H(o_0)
\ \middle|\
c,a,h_v
\right].
\label{eqoptioneffect}
\end{equation}
Here causal refers to the effect of invoking a search tactic relative to its matched continuation. It does not refer to causal laws inside the scientific domain. Known assignment probabilities and branch isolation allow the registry to estimate this effect within the audited context. Harnesses can be compared in the same way by changing the harness version while holding the state, role, option, horizon, cost, and verifier fixed.

An option enters long-term causal memory after repeated audit evidence supports a positive contextual effect. Retrieval is conditioned on the current state, role, harness, and horizon. A synthesized natural-language strategy is evaluated as a new option before it receives the status of an established memory. A small no-memory audit arm continues to measure whether retrieval remains useful over time.

Feedback follows three clear paths. Verified outcomes update the archive and scheduler. Matched audit outcomes update causal memory. Homogeneous groups of verified branches update the role adapters. Keeping these paths separate makes it clear whether an observation supports search allocation, causal attribution, or policy learning.

\subsection{Test-Time Policy Learning}
\label{seclearning}

This is the reinforcement-learning component of EVOLVE. Learning occurs after a group of branches has closed. Every branch in a learning group shares the same scientific cell, role, option, harness version, horizon, cost range, policy snapshot, and frozen record threshold. This makes the reward comparison meaningful.

Suppose $K$ branches are sampled from one frozen role policy and their verified gains are sorted as $G_{(1,K)}\leq\cdots\leq G_{(K,K)}$. EVOLVE optimizes the mean of the largest $m$ gains
\begin{equation}
J_{K,m}(\theta)
=
\mathbb E\left[
\frac{1}{m}
\sum_{j=K-m+1}^{K}G_{(j,K)}
\right].
\label{eqtopmk}
\end{equation}
Pure Max at $K$ is recovered when $m=1$. A slightly larger $m$ keeps learning focused on the upper tail while reducing sensitivity to a single extreme outcome.

A larger on-policy batch with $N>K$ branches can be reused across subsets of size $K$. OrderGrad provides the corresponding finite-sample policy-gradient estimator under its sampling assumptions \citep{parmas2026}. MaxPO provides an exactly centered alternative for the pure maximum objective \citep{takashiro2026}. The log probability of a branch is the sum of the role policy log probabilities for its internal decisions. After a group closes, its gradient updates the LoRA adapter of that role. The shared backbone and the other two role adapters remain unchanged.

The scheduler and the role policies learn different things. The scheduler learns which allocation is likely to produce the next record. Role learning improves how a selected option is carried out. Both use verified descendant maxima, which keeps search and training aligned with Eq.~\eqref{eqobjective}.

\subsection{Bounded Refinement}
\label{secrefinement}

A provisional candidate with a clear and actionable diagnosis may enter a bounded refinement stage. It receives one to three attempts, a fixed cost limit, a maximum depth of two, and no second entry after leaving the stage. The Challenger proposes the smallest change that may address the diagnosis. Each revision is treated as a new candidate and returns to the blinded verifier with a new evidence packet.

Refinement is also represented as an option. Eligible provisional candidates in the audit channel are randomly assigned either to refinement or to an equal-cost fresh continuation from the same starting state. This comparison estimates when refinement is useful. A refined candidate updates the archive, causal memory, or role learning data after independent confirmation.

When the total budget is exhausted, EVOLVE returns the verified record holder together with its evidence packet and complete lineage.

\clearpage

\subsection{Detailed Framework Views}
\label{secdetailedviews}

The following diagrams expand the main EVOLVE cycle into six complementary views. Each view follows the same evidence flow while isolating one part of the framework that benefits from a closer explanation.

\paragraph{End-to-End Execution}
The complete pipeline begins from the verified scientific archive, promoted causal option records, and versioned harness pool. The scheduler forms allocation arms, protects support for audits and coverage, and dispatches production branches and randomized audit pairs through isolated role policies. Independent verification routes evidence to the archive, scheduler, causal memory, and role learning before accepted updates cross the synchronization barrier together.

\begin{figure}[!htbp]
\centering
\resizebox{\linewidth}{!}{\begin{tikzpicture}[
  x=1cm,
  y=1cm,
  font=\sffamily\scriptsize,
  box/.style={draw, rounded corners=2pt, thick, align=center, minimum height=8mm, fill=white},
  state/.style={box, draw=blue!65!black, fill=blue!6},
  control/.style={box, draw=violet!65!black, fill=violet!7},
  role/.style={box, draw=teal!65!black, fill=teal!7},
  branch/.style={box, draw=orange!75!black, fill=orange!8},
  verify/.style={box, draw=green!55!black, fill=green!8},
  memory/.style={box, draw=green!45!black, fill=green!5},
  update/.style={box, draw=cyan!55!black, fill=cyan!6},
  refine/.style={box, draw=orange!85!black, fill=orange!13},
  neutral/.style={box, draw=black!55, fill=black!3},
  flow/.style={-{Latex[length=2mm]}, thick},
  auditflow/.style={-{Latex[length=2mm]}, thick, draw=violet!75!black},
  nextflow/.style={-{Latex[length=2mm]}, thick, dashed, draw=black!60},
  smalllabel/.style={font=\sffamily\tiny, align=center, text=black!75}
]

\node[state, minimum width=35mm] (archive) at (-4.8,0) {Verified scientific archive\\quality states \quad stepping stones};
\node[memory, minimum width=35mm] (registry) at (0,0) {Promoted causal option records\\contextual effect evidence};
\node[control, minimum width=35mm] (harnesses) at (4.8,0) {Versioned harness pool\\stable and trial versions};

\node[control, minimum width=83mm] (builder) at (0,-1.55) {Allocation arm builder\\$r=(c,a,o,h_v,H,C)$};
\draw[flow] (archive) -- (builder);
\draw[flow] (registry) -- (builder);
\draw[flow] (harnesses) -- (builder);

\node[control, minimum width=103mm, minimum height=14mm] (scheduler) at (0,-3.15) {Budgeted posterior controller\\audit support \quad cell coverage\\role and harness support \quad record-improvement portfolio};
\draw[flow] (builder) -- (scheduler);

\node[role, minimum width=33mm] (scout) at (-4.2,-4.75) {Scout LoRA branch\\new regions and options};
\node[role, minimum width=33mm] (mechanist) at (0,-4.75) {Mechanist LoRA branch\\mechanisms and constraints};
\node[role, minimum width=33mm] (challenger) at (4.2,-4.75) {Challenger LoRA branch\\counterexamples and tests};
\draw[flow] (scheduler) -- (scout);
\draw[flow] (scheduler) -- (mechanist);
\draw[flow] (scheduler) -- (challenger);

\node[neutral, minimum width=64mm, minimum height=10mm] (dispatch) at (0,-5.95) {Role-isolated branch executor\\one frozen role checkpoint per branch};
\draw[flow] (scout) -- (dispatch);
\draw[flow] (mechanist) -- (dispatch);
\draw[flow] (challenger) -- (dispatch);

\node[branch, minimum width=39mm] (production) at (-3.1,-7.15) {Production branches\\seek a new verified record};
\node[branch, minimum width=50mm, draw=violet!70!black, fill=violet!6] (audit) at (3.1,-7.15) {Randomized audit pairs\\tested option \quad matched continuation};
\draw[flow] (dispatch) -- (production);
\draw[auditflow] (dispatch) -- (audit);

\node[verify, minimum width=90mm, minimum height=15mm] (verifier) at (0,-8.75) {Common independent verifier\\admissibility \quad scientific measure \quad perturbations\\replication \quad blinded checks};
\draw[flow] (production) -- node[smalllabel, left] {closed production evidence} (verifier);
\draw[auditflow] (audit) -- node[smalllabel, right] {paired audit evidence} (verifier);

\node[state, minimum width=27mm] (archiveupdate) at (-5.1,-10.75) {Archive and\\record update};
\node[control, minimum width=27mm] (posterior) at (-1.7,-10.75) {Tail posterior\\update};
\node[memory, minimum width=27mm] (memoryupdate) at (1.7,-10.75) {Causal option\\memory update};
\node[update, minimum width=27mm] (groups) at (5.1,-10.75) {Homogeneous\\learning groups};

\draw[flow] (verifier.south west) --
  node[smalllabel, pos=.42, below, sloped, inner sep=.7pt]
       {admitted\\states}
  (archiveupdate.north east);
\draw[flow] (verifier.south) -- node[smalllabel, pos=.58, above, sloped, fill=white, inner sep=1pt] {production} (posterior.north);
\draw[auditflow] (verifier.south) -- node[smalllabel, pos=.58, above, sloped, fill=white, inner sep=1pt] {paired effects} (memoryupdate.north);
\draw[flow] (verifier.south east) -- node[smalllabel, pos=.62, above, sloped, fill=white, inner sep=1pt] {role data} (groups.north west);

\node[refine, minimum width=42mm] (provisional) at (8.15,-8.55) {Provisional candidate\\with actionable diagnosis};
\node[refine, minimum width=42mm] (refinement) at (8.15,-9.70) {Matched nursery allocation\\refinement versus fresh continuation};
\draw[flow, draw=orange!80!black] (verifier.east) -- (provisional.west);
\draw[flow, draw=orange!80!black] (provisional) -- (refinement);
\coordinate (reverifyout) at ([xshift=-1mm]refinement.west);
\draw[flow, draw=orange!80!black]
  (refinement.west) -- (reverifyout)
  |- ([yshift=-3.2mm]verifier.east);
\node[smalllabel, anchor=north east, fill=white, inner sep=.7pt]
  at (5.80,-9.16) {blinded\\re-verification};

\node[update, minimum width=42mm] (loraupdate) at (5.1,-12.15) {Role-specific LoRA update\\shared backbone remains frozen};
\draw[flow] (groups) -- (loraupdate);

\node[neutral, minimum width=111mm, minimum height=10mm] (barrier) at (0,-13.6) {Synchronization barrier\\verified updates enter the next epoch};
\draw[nextflow] (archiveupdate) -- (barrier);
\draw[nextflow] (posterior) -- (barrier);
\draw[nextflow] (memoryupdate) -- (barrier);
\draw[nextflow] (loraupdate) -- (barrier);

\node[font=\sffamily\tiny\bfseries, text=black!75, anchor=west]
  at (-6.55,0.72) {Inputs available at the start of an epoch};
\node[font=\sffamily\tiny\bfseries, text=black!75, anchor=west]
  at (-5.85,-4.15) {Parallel role execution};
\end{tikzpicture}
}
\caption{Detailed EVOLVE execution across allocation, isolated role branches, independent verification, bounded refinement, and synchronized updates.}
\label{figdetailedpipeline}
\end{figure}

\clearpage

\paragraph{Archive and Provenance}
The scientific archive is persistent, while every launched option creates a local tree of descendants. The verifier and frozen descriptor map decide whether a candidate is admitted and which scientific cell receives it. A separate candidate-level provenance graph preserves route identity and complete lineage even when different branches reach the same cell.

\begin{figure}[!htbp]
\centering
\resizebox{\linewidth}{!}{\begin{tikzpicture}[
  x=1cm,
  y=1cm,
  font=\sffamily\scriptsize,
  cell/.style={draw=blue!65!black, thick, rounded corners=1.5pt, fill=blue!5, minimum width=20mm, minimum height=14mm, align=center},
  selected/.style={cell, fill=blue!13, very thick},
  state/.style={circle, draw=blue!70!black, thick, fill=white, minimum size=5.5mm, inner sep=0pt, font=\sffamily\tiny},
  launch/.style={draw=orange!80!black, thick, rounded corners=2pt, fill=orange!10, minimum width=18mm, minimum height=7mm, align=center},
  intermediate/.style={circle, draw=black!55, thick, fill=black!4, minimum size=5mm, inner sep=0pt},
  admitted/.style={circle, draw=green!55!black, thick, fill=green!12, minimum size=6mm, inner sep=0pt, font=\sffamily\tiny},
  rejected/.style={circle, draw=black!45, thick, fill=black!8, minimum size=6mm, inner sep=0pt, font=\sffamily\tiny},
  process/.style={draw, thick, rounded corners=2pt, fill=white, minimum width=25mm, minimum height=9mm, align=center},
  flow/.style={-{Latex[length=2mm]}, thick},
  execution/.style={-{Latex[length=2mm]}, thick, draw=orange!80!black},
  mapping/.style={-{Latex[length=2mm]}, thick, draw=blue!70!black},
  provenance/.style={-{Latex[length=2mm]}, thick, draw=violet!70!black},
  note/.style={font=\sffamily\tiny, align=center, text=black!70}
]

\node[font=\sffamily\small\bfseries, anchor=west] at (-6.5,1.35) {Frozen archive at epoch start};

\node[cell] (cb) at (-5.2,0) {Under-tested\\cell $c_B$};
\node[selected] (ca) at (-3.0,0) {};
\node[note, text=blue!70!black] at (-3.0,0.56) {Cell $c_A$};
\node[state] (qa) at (-3.60,-0.05) {Q};
\node[state] (va) at (-3.00,-0.36) {V};
\node[state] (sa) at (-2.40,-0.05) {S};
\node[note, text width=50mm] at (-3.0,-0.90)
  {$Q$ quality \quad $V$ descendant value \quad $S$ stepping stone};

\node[cell, dashed] (cc) at (-5.2,-1.9) {Empty\\cell $c_C$};
\node[cell] (cd) at (-3.0,-1.9) {Cell $c_D$\\stepping stone};
\node[process, draw=yellow!60!black, fill=yellow!13, minimum width=30mm] (record) at (-4.1,-3.35) {Protected global record\\$M_e^{\mathrm{start}}$};

\node[font=\sffamily\small\bfseries] at (0.45,1.35) {Local option trees};
\node[launch, minimum height=9mm] (l1) at (-0.55,0.25) {Scout\\option $o_1$};
\node[intermediate] (y1) at (1.10,0.25) {};
\node[admitted] (x1) at (2.15,0.95) {$x_1$};
\node[admitted] (x2) at (2.15,0.20) {$x_2$};
\node[rejected] (x3) at (2.15,-0.55) {$\times$};

\node[launch, minimum height=9mm] (l2) at (-0.55,-1.8) {Mechanist\\option $o_2$};
\node[intermediate] (y2) at (1.10,-1.8) {};
\node[admitted] (x4) at (2.15,-1.45) {$x_4$};
\node[rejected] (x5) at (2.15,-2.20) {$\times$};

\coordinate (qaout) at (-1.90,0.30);
\draw[execution] (qa.north) |- (qaout) -- (l1.west);
\draw[execution] (l1) -- (y1);
\draw[execution] (y1) -- (x1);
\draw[execution] (y1) -- (x2);
\draw[execution] (l1.south east) -- (x3.west);
\draw[execution] (cd.east) -- (l2.west);
\draw[execution] (l2) -- (y2);
\draw[execution] (y2) -- (x4);
\draw[execution] (y2) -- (x5);

\node[note, draw=orange!70!black, rounded corners=1pt, fill=orange!5, text width=50mm] at (0.2,-3.25) {$Z_H$ is the best verified descendant\\$G_H=[Z_H-M_e^{\mathrm{start}}]_+$};

\draw[dashed, thick, black!50] (3.15,1.15) -- (3.15,-3.70);
\node[note, rotate=90, fill=white, inner sep=1pt] at (3.15,-2.55) {verification and synchronization};

\node[process, draw=green!55!black, fill=green!7, minimum width=28mm] (verifier) at (4.90,0.15) {Independent verifier};
\node[process, draw=blue!65!black, fill=blue!6] (mapper) at (4.90,-1.15) {Descriptor and cell map\\$q_v\!\circ\!\phi_v$};

\draw[thick, black!55] (2.80,0.95) -- (2.80,-2.20);
\draw[dashed, thick, black!55] (x1.east) -- (2.80,0.95);
\draw[dashed, thick, black!55] (x2.east) -- (2.80,0.20);
\draw[dashed, thick, black!55] (x3.east) -- (2.80,-0.55);
\draw[dashed, thick, black!55] (x4.east) -- (2.80,-1.45);
\draw[dashed, thick, black!55] (x5.east) -- (2.80,-2.20);
\draw[flow] (2.80,0.15) -- (verifier.west);
\fill[black!55] (2.80,0.15) circle (1pt);
\draw[mapping] (verifier) -- (mapper);

\node[font=\sffamily\small\bfseries] at (9.20,1.35) {Updated archive and provenance};
\node[cell, minimum height=16mm] (newcc) at (7.90,-0.10) {};
\node[note, text=blue!70!black] at (7.90,0.38) {Cell $c_C$};
\node[cell, minimum width=26mm, minimum height=16mm] (newcb) at (10.55,-0.10) {};
\node[note, text=blue!70!black] at (10.55,0.38) {Cell $c_B$};
\node[admitted] (newx1) at (7.90,-0.31) {$x_1$};
\node[admitted] (newx2) at (10.20,-0.31) {$x_2$};
\node[admitted] (newx4) at (10.90,-0.31) {$x_4$};
\node[process, draw=yellow!60!black, fill=yellow!13, minimum width=26mm] (newrecord) at (10.55,-1.55) {New global record\\$x_2$};

\coordinate (mapfork) at (6.55,-0.10);
\draw[mapping, -] (mapper.east) -- (6.55,-1.15) -- (mapfork);
\draw[mapping] (mapfork) -- (newcc.west);
\draw[mapping] (mapfork)
  -- (6.55,1.05) -- (10.55,1.05) -- (newcb.north);
\fill[blue!70!black] (mapfork) circle (1pt);
\coordinate (recordexit) at (newx2.south |- newcb.south);
\draw[flow, draw=yellow!65!black]
  (newx2.south) -- (recordexit)
  -- ([xshift=-3.5mm]newrecord.north);

\node[process, draw=violet!65!black, fill=violet!5, minimum width=60mm, minimum height=20mm] (graphbox) at (8.50,-3.30) {};
\node[note, text=violet!70!black] at (8.50,-2.55) {Persistent candidate-level provenance graph};
\node[state] (pqa) at (6.00,-3.40) {Q};
\node[intermediate] (py1) at (7.00,-3.40) {};
\node[admitted] (px1) at (8.10,-3.05) {$x_1$};
\node[admitted] (px2) at (8.10,-3.75) {$x_2$};
\node[state] (psd) at (8.90,-3.40) {S};
\node[intermediate] (py2) at (9.75,-3.40) {};
\node[admitted] (px4) at (10.60,-3.40) {$x_4$};
\draw[provenance] (pqa) -- (py1);
\draw[provenance] (py1) -- (px1);
\draw[provenance] (py1) -- (px2);
\draw[provenance] (psd) -- (py2);
\draw[provenance] (py2) -- (px4);

\node[note] at (8.50,-4.60) {Candidates may share a scientific cell while retaining distinct identities and lineages};

\end{tikzpicture}
}
\caption{Archive cells, local option trees, descriptor-based placement, and the persistent candidate provenance graph.}
\label{figarchiveprovenance}
\end{figure}

\clearpage

\paragraph{Agent and Memory Isolation}
Scout, Mechanist, and Challenger share one read-only backbone and a verified public evidence layer. Each role keeps its own LoRA adapter, optimizer, random stream, working memory, retrieval view, and raw branch transcript. Only compact verified packets cross the synchronization barrier, and each role learns from its own homogeneous verified group.

\begin{figure}[!htbp]
\centering
\resizebox{\linewidth}{!}{\begin{tikzpicture}[
  x=1cm,
  y=1cm,
  font=\sffamily\scriptsize,
  shared/.style={draw=black!60, thick, rounded corners=2pt, fill=black!4, minimum height=8mm, align=center},
  rolebox/.style={draw=teal!65!black, thick, rounded corners=2pt, fill=teal!7, minimum width=34mm, minimum height=8mm, align=center},
  private/.style={draw=teal!55!black, thick, rounded corners=2pt, fill=white, minimum width=34mm, minimum height=7mm, align=center},
  verify/.style={draw=green!55!black, thick, rounded corners=2pt, fill=green!7, minimum height=9mm, align=center},
  update/.style={draw=cyan!55!black, thick, rounded corners=2pt, fill=cyan!7, minimum width=34mm, minimum height=9mm, align=center},
  barrier/.style={draw=violet!60!black, thick, rounded corners=2pt, fill=violet!6, minimum height=8mm, align=center},
  flow/.style={-{Latex[length=2mm]}, thick},
  evidence/.style={-{Latex[length=2mm]}, thick, draw=green!55!black},
  note/.style={font=\sffamily\tiny, align=center, text=black!70}
]

\node[shared, minimum width=115mm] (backbone) at (0,0) {Frozen foundation model \quad read only for the full discovery run};
\node[shared, minimum width=115mm] (public) at (0,-1.15) {Verified shared layer \quad scientific archive \quad promoted causal options \quad versioned harnesses};

\node[font=\sffamily\small\bfseries, text=teal!70!black] at (-4.0,-2.25) {Scout};
\node[font=\sffamily\small\bfseries, text=teal!70!black] at (0,-2.25) {Mechanist};
\node[font=\sffamily\small\bfseries, text=teal!70!black] at (4.0,-2.25) {Challenger};

\node[rolebox] (slora) at (-4.0,-3.05) {Scout LoRA\\private optimizer};
\node[rolebox] (mlora) at (0,-3.05) {Mechanist LoRA\\private optimizer};
\node[rolebox] (clora) at (4.0,-3.05) {Challenger LoRA\\private optimizer};

\node[private] (smemory) at (-4.0,-4.15) {Private working memory\\and random stream};
\node[private] (mmemory) at (0,-4.15) {Private working memory\\and random stream};
\node[private] (cmemory) at (4.0,-4.15) {Private working memory\\and random stream};

\node[private] (sretrieval) at (-4.0,-5.25) {Role retrieval view\\verified packets only};
\node[private] (mretrieval) at (0,-5.25) {Role retrieval view\\verified packets only};
\node[private] (cretrieval) at (4.0,-5.25) {Role retrieval view\\verified packets only};

\node[rolebox] (sbranch) at (-4.0,-6.4) {Scout option branches\\raw transcript stays here};
\node[rolebox] (mbranch) at (0,-6.4) {Mechanist option branches\\raw transcript stays here};
\node[rolebox] (cbranch) at (4.0,-6.4) {Challenger option branches\\raw transcript stays here};

\draw[flow] (slora) -- (smemory);
\draw[flow] (smemory) -- (sretrieval);
\draw[flow] (sretrieval) -- (sbranch);
\draw[flow] (mlora) -- (mmemory);
\draw[flow] (mmemory) -- (mretrieval);
\draw[flow] (mretrieval) -- (mbranch);
\draw[flow] (clora) -- (cmemory);
\draw[flow] (cmemory) -- (cretrieval);
\draw[flow] (cretrieval) -- (cbranch);

\node[verify, minimum width=115mm] (verifier) at (0,-7.85) {Common independent verifier \quad same admission rule and scientific measure for every role};
\draw[flow] (sbranch.south) -- ([xshift=-40mm]verifier.north);
\draw[flow] (mbranch.south) -- (verifier.north);
\draw[flow] (cbranch.south) -- ([xshift=40mm]verifier.north);

\node[barrier, minimum width=115mm] (barrier) at (0,-9.05) {Synchronization barrier \quad immutable verified packets cross role boundaries};
\draw[evidence] (verifier) -- (barrier);

\node[update] (supdate) at (-4.0,-10.35) {Scout verified group\\updates Scout LoRA};
\node[update] (mupdate) at (0,-10.35) {Mechanist verified group\\updates Mechanist LoRA};
\node[update] (cupdate) at (4.0,-10.35) {Challenger verified group\\updates Challenger LoRA};
\draw[evidence] ([xshift=-40mm]barrier.south) -- (supdate.north);
\draw[evidence] (barrier.south) -- (mupdate.north);
\draw[evidence] ([xshift=40mm]barrier.south) -- (cupdate.north);

\node[note, text width=105mm] at (0,-11.25) {The archive and promoted causal records become public in the next epoch\\Role gradients remain separate};

\end{tikzpicture}
}
\caption{Isolation of role adapters, working memories, retrieval views, branch transcripts, and verified learning updates.}
\label{figagentisolation}
\end{figure}

\clearpage

\paragraph{Adaptive Harness Allocation}
The harness version is part of every allocation arm alongside the state, role, option, horizon, and cost. The scheduler reserves support for audits, scientific coverage, every role, and provisional harness trials before allocating the remaining budget toward expected record improvement. A selected harness remains locked throughout its branch, while later evidence can change which version is chosen for future branches.

\begin{figure}[!htbp]
\centering
\resizebox{\linewidth}{!}{\begin{tikzpicture}[
  x=1cm,
  y=1cm,
  font=\sffamily\scriptsize,
  box/.style={draw, rounded corners=2pt, thick, align=center, minimum height=9mm, inner sep=3pt, fill=white},
  state/.style={box, draw=blue!65!black, fill=blue!6},
  harness/.style={box, draw=violet!65!black, fill=violet!7},
  trial/.style={harness, dashed},
  control/.style={box, draw=violet!70!black, fill=violet!5},
  branch/.style={box, draw=orange!75!black, fill=orange!8},
  verify/.style={box, draw=green!55!black, fill=green!8},
  evidence/.style={box, draw=green!55!black, fill=green!5},
  neutral/.style={box, draw=black!55, fill=black!3},
  gate/.style={circle, draw=violet!70!black, thick, fill=violet!7, minimum size=8mm, inner sep=1pt},
  flow/.style={-{Latex[length=2mm]}, thick},
  auditflow/.style={-{Latex[length=2mm]}, thick, draw=violet!75!black},
  nextflow/.style={-{Latex[length=2mm]}, thick, dashed, draw=black!60},
  note/.style={font=\sffamily\tiny, align=center, text=black!70}
]

\node[state, text width=33mm] (context) at (-7.60,0.80) {Scientific context\\cells \quad roles \quad options\\horizon and cost menu};
\node[harness, text width=33mm] (registry) at (-7.60,-0.90) {Versioned harness registry\\stable $h_1^{(3)}$ \quad stable $h_2^{(5)}$\\trial $h_3^{(1)}$};

\node[control, text width=35mm, minimum height=15mm] (builder) at (-3.60,0) {Arm constructor\\$r=(c,a,o,h_v,H,C)$};
\draw[flow] (context.east) -- ([yshift=4mm]builder.west);
\draw[flow] (registry.east) -- ([yshift=-4mm]builder.west);

\node[control, text width=42mm, minimum height=20mm] (posterior) at (1.20,0) {Joint outcome posterior\\admission probability\\admitted upper tail\\dependence between arms};
\draw[flow] (builder.east) -- (posterior.west);

\node[neutral, draw=yellow!65!black, fill=yellow!12, text width=34mm] (record) at (4.00,2.00) {Current verified record $M_t$};
\node[neutral, text width=39mm] (quotas) at (8.10,2.00) {Reserved support\\audit \quad coverage\\roles \quad harness trials};
\node[control, text width=43mm, minimum height=20mm] (optimizer) at (6.10,0) {Budgeted portfolio selection\\maximize posterior record improvement\\under remaining cost};
\draw[flow] (posterior.east) -- (optimizer.west);
\draw[flow] (record.south) -- ([xshift=-7mm]optimizer.north);
\draw[flow] (quotas.south) -- ([xshift=7mm]optimizer.north);

\node[branch, text width=38mm, minimum height=25mm] (portfolio) at (-3.60,-2.80) {Selected portfolio $S_t$\\$r_1$ with $h_1^{(3)}$ locked\\$r_2$ with $h_2^{(5)}$ locked\\$r_3$ with trial $h_3^{(1)}$ locked};
\coordinate (selectrail) at ([yshift=-2.5mm]optimizer.south);
\coordinate (selectdrop) at (portfolio.north |- selectrail);
\draw[flow] (optimizer.south) -- (selectrail)
  -- node[note, pos=.20, below=1pt, fill=white, inner sep=1pt] {selected arms}
  (selectdrop) -- (portfolio.north);

\node[branch, text width=32mm, minimum height=15mm] (execution) at (0.80,-2.80) {Locked branch execution\\one version for the branch lifetime};
\node[verify, text width=34mm, minimum height=20mm] (verifier) at (4.70,-2.80) {Fixed external verifier\\same admission rule\\for every harness};
\draw[flow] (portfolio.east) -- (execution.west);
\draw[flow] (execution.east) -- (verifier.west);

\node[evidence, text width=36mm] (production) at (2.20,-4.75) {Production evidence\\absolute record value};
\node[evidence, draw=violet!65!black, fill=violet!6, text width=36mm] (audit) at (7.00,-4.75) {Matched harness audit\\contextual difference};
\draw[flow] ([xshift=-5mm]verifier.south) -- ++(0,-0.15) -| (production.north);
\draw[auditflow] ([xshift=5mm]verifier.south) -- ++(0,-0.15) -| (audit.north);

\node[control, text width=34mm] (nextposterior) at (2.20,-6.35) {Tail posterior update\\next epoch};
\node[harness, text width=34mm] (effect) at (7.00,-6.35) {Harness effect record\\next epoch};
\draw[nextflow] (production.south) -- (nextposterior.north);
\draw[nextflow] (audit.south) -- (effect.north);

\node[gate] (versiongate) at (4.60,-5.90) {$\Delta h$};
\node[trial, text width=36mm] (newversion) at (4.60,-7.45) {New trial version $h_3^{(2)}$\\joins the next epoch};
\draw[nextflow] (production.south east) -- (versiongate.north west);
\draw[nextflow] (audit.south west) -- (versiongate.north east);
\draw[nextflow] (versiongate.south) -- (newversion.north);

\node[neutral, text width=170mm, minimum height=14mm] (rule) at (1.00,-8.90) {Adaptation occurs only between branches\\Harness versions remain locked within every active branch\\All comparisons use the same external verifier};

\end{tikzpicture}
}
\caption{Adaptive allocation over state, role, option, harness version, horizon, and cost with locked branch execution.}
\label{figharnessallocation}
\end{figure}

\clearpage

\paragraph{Randomized Audits and Causal Memory}
An option branch and its matched continuation begin from the same frozen state and role checkpoint under the same harness, horizon, cost, and verifier. Their paired difference estimates the local effect of invoking the option rather than the selection preference of the scheduler. Repeated positive audit evidence can promote a context-conditioned record into causal option memory, while uncertain evidence leads to further matched audits.

\begin{figure}[!htbp]
\centering
\resizebox{\linewidth}{!}{\begin{tikzpicture}[
  x=1cm,
  y=1cm,
  font=\sffamily\scriptsize,
  box/.style={draw, rounded corners=2pt, thick, align=center, minimum height=8mm, inner sep=3pt, fill=white},
  frozen/.style={box, draw=black!60, fill=black!4},
  audit/.style={box, draw=violet!70!black, fill=violet!7},
  verify/.style={box, draw=green!55!black, fill=green!7},
  evidence/.style={box, draw=green!50!black, fill=green!5},
  memory/.style={box, draw=blue!60!black, fill=blue!6},
  flow/.style={-{Latex[length=2mm]}, thick},
  auditflow/.style={-{Latex[length=2mm]}, thick, draw=violet!75!black},
  promotion/.style={-{Latex[length=2mm]}, thick, dashed, draw=green!55!black},
  note/.style={font=\sffamily\tiny, align=center, text=black!70}
]

\node[font=\sffamily\small\bfseries, anchor=west] at (-8.2,1.55) {A \quad Matched randomized audit};

\node[frozen, text width=31mm, minimum height=14mm] (start) at (-6.60,-0.25) {Frozen start packet\\same cell and role checkpoint};
\node[note, anchor=north] at ([yshift=-2mm]start.south) {same harness \quad horizon \quad cost \quad verifier};

\node[audit, text width=25mm] (test) at (-3.00,0.35) {Test branch\\option $o$};
\node[audit, text width=25mm] (control) at (-3.00,-0.85) {Control branch\\continuation $o_0$};
\coordinate (assignmentfork) at (-4.65,-0.25);
\draw[thick, draw=violet!75!black] (start.east) -- (assignmentfork);
\fill[violet!75!black] (assignmentfork) circle (1pt);
\draw[auditflow] (assignmentfork) |- (test.west);
\draw[auditflow] (assignmentfork) |- (control.west);
\node[note, fill=white, inner sep=1.5pt] at (-4.65,1.02) {random assignment};

\node[verify, text width=20mm] (vtest) at (0,0.35) {Same verifier};
\node[verify, text width=20mm] (vcontrol) at (0,-0.85) {Same verifier};
\node[evidence, text width=20mm] (gtest) at (2.80,0.35) {Gain $G_H(o)$};
\node[evidence, text width=20mm] (gcontrol) at (2.80,-0.85) {Gain $G_H(o_0)$};
\draw[auditflow] (test) -- (vtest);
\draw[auditflow] (control) -- (vcontrol);
\draw[flow] (vtest) -- (gtest);
\draw[flow] (vcontrol) -- (gcontrol);

\node[audit, text width=35mm, minimum height=13mm] (effect) at (6.30,-0.25) {Paired branch effect\\$\Delta G=G_H(o)-G_H(o_0)$};
\coordinate (effectmerge) at (4.05,-0.25);
\draw[thick, draw=violet!75!black] (gtest.east) -- (4.05,0.35) -- (effectmerge);
\draw[thick, draw=violet!75!black] (gcontrol.east) -- (4.05,-0.85) -- (effectmerge);
\fill[violet!75!black] (effectmerge) circle (1pt);
\draw[auditflow] (effectmerge) -- (effect.west);

\draw[dashed, thick, black!45] (-8.2,-2.05) -- (8.2,-2.05);
\node[font=\sffamily\small\bfseries, anchor=west] at (-8.2,-2.65) {B \quad Evidence promotion across repeated audit pairs};

\node[evidence, text width=29mm] (blocks) at (-5.90,-3.75) {Closed audit blocks};
\node[evidence, text width=39mm] (estimate) at (-1.80,-3.75) {Contextual effect estimate\\and confidence sequence};
\node[memory, text width=34mm, minimum height=12mm] (decision) at (2.60,-3.75) {Repeated positive\\lower bound};
\draw[flow] (blocks) -- (estimate);
\draw[flow] (estimate) -- (decision);

\node[frozen, text width=30mm] (more) at (0.20,-5.30) {Provisional evidence\\collect more audits};
\node[memory, text width=32mm] (promoted) at (4.60,-5.30) {Promoted causal\\option record};
\coordinate (evidencefork) at (2.60,-4.55);
\draw[flow, -] (decision.south) -- (evidencefork);
\draw[flow] (evidencefork) -| (more.north);
\node[note, fill=white, inner sep=1pt] at (1.40,-4.70) {not yet};
\draw[promotion] (evidencefork) -| (promoted.north);
\node[note, fill=white, inner sep=1pt] at (3.60,-4.70) {supported};
\fill[black!70] (evidencefork) circle (1pt);

\node[memory, text width=39mm] (retrieval) at (4.60,-6.75) {Context-conditioned retrieval\\state \quad role \quad harness \quad horizon};
\node[audit, text width=33mm, dashed] (nomemory) at (0.20,-6.75) {Randomized no-memory\\audit remains active};
\draw[promotion] (promoted) -- (retrieval);
\draw[auditflow] (more) -- (nomemory);

\node[frozen, text width=155mm, minimum height=9mm] at (0,-8.15) {Production branches update scheduling \qquad randomized matched comparisons update causal option memory};

\node[note,
  draw=violet!50!black,
  rounded corners=1pt,
  fill=white,
  text width=30mm,
  anchor=south west]
  at ([xshift=2mm,yshift=1mm]decision.north east)
  {Each effect remains local\\to its audited context};

\end{tikzpicture}
}
\caption{Matched randomized audit pairs and evidence promotion into context-conditioned causal option memory.}
\label{figauditmemory}
\end{figure}

\clearpage

\paragraph{Test-Time Learning and Refinement}
Homogeneous on-policy groups produce an upper-tail update for one selected role while the backbone and the other role adapters remain fixed. Provisional outcomes follow a separate bounded path that permits only a few diagnosis-guided revisions. Every revision returns to independent verification, and only confirmed revisions enter their own refinement learning group.

\begin{figure}[!htbp]
\centering
\resizebox{\linewidth}{!}{\begin{tikzpicture}[
  x=1cm,
  y=1cm,
  font=\sffamily\scriptsize,
  box/.style={draw, rounded corners=2pt, thick, align=center, minimum height=9mm, inner sep=3pt, fill=white},
  frozen/.style={box, draw=black!60, fill=black!4},
  role/.style={box, draw=teal!65!black, fill=teal!7},
  branch/.style={box, draw=orange!75!black, fill=orange!8},
  verify/.style={box, draw=green!55!black, fill=green!7},
  gain/.style={box, draw=blue!65!black, fill=blue!6},
  learning/.style={box, draw=cyan!60!black, fill=cyan!7},
  refine/.style={box, draw=orange!85!black, fill=orange!13},
  flow/.style={-{Latex[length=2mm]}, thick},
  refineflow/.style={-{Latex[length=2mm]}, thick, draw=orange!80!black},
  separate/.style={-{Latex[length=2mm]}, thick, dashed, draw=orange!80!black},
  note/.style={font=\sffamily\tiny, align=center, text=black!70}
]

\node[frozen, text width=143mm, minimum height=12mm] at (0,1.20) {\textbf{One homogeneous group}\\same cell \quad role \quad option \quad harness \quad horizon \quad cost \quad checkpoint \quad record threshold};

\node[role, text width=36mm] (snapshot) at (-5.1,0) {Frozen selected-role LoRA\\$\theta_a^e$};
\node[branch, text width=36mm] (branches) at (0,0) {$N$ on-policy option branches\\from one frozen snapshot};
\node[verify, text width=36mm] (verifier) at (5.1,0) {Independent verification\\completed branches only};
\draw[flow] (snapshot) -- (branches);
\draw[flow] (branches) -- (verifier);

\node[gain, text width=36mm] (gains) at (5.1,-1.75) {Confirmed descendant gains\\$G_1,\ldots,G_N$};
\node[learning, text width=40mm] (credit) at (0,-1.75) {Top-$m$-at-$K$ credit\\OrderGrad \quad MaxPO when $m=1$};
\node[role, text width=38mm] (update) at (-5.1,-1.75) {Update selected-role LoRA\\for epoch $e+1$};
\draw[flow] (verifier.south) -- (gains.north);
\draw[flow] (gains.west) -- (credit.east);
\draw[flow] (credit.west) -- (update.east);

\node[frozen, text width=143mm, minimum height=12mm] at (0,-3.1) {\textbf{Only the selected-role LoRA updates}\\shared backbone fixed \quad other role LoRAs unchanged \quad no cross-role gradient \quad update after the group closes};

\draw[dashed, thick, black!45] (-7.8,-3.90) -- (7.8,-3.90);
\node[font=\sffamily\small\bfseries, anchor=west, text=orange!75!black] at (-7.2,-4.35) {Separate path for a provisional verifier outcome};

\node[refine, text width=24mm, minimum height=13mm, inner sep=2pt] (diagnosis) at (-6.4,-5.4) {Actionable\\diagnosis};
\node[refine, text width=24mm, minimum height=13mm, inner sep=2pt] (refinement) at (-3.2,-5.4) {Bounded refinement\\option\\1--3 attempts};
\node[verify, text width=24mm, minimum height=13mm, inner sep=2pt] (reverify) at (0,-5.4) {Blinded\\re-verification};
\node[gain, text width=24mm, minimum height=13mm, inner sep=2pt] (confirmed) at (3.2,-5.4) {Confirmed revised\\candidate};
\node[refine, text width=24mm, minimum height=13mm, inner sep=2pt] (refinegroup) at (6.4,-5.4) {Separate refinement\\learning group};
\draw[refineflow] (diagnosis) -- (refinement);
\draw[refineflow] (refinement) -- (reverify);
\draw[refineflow] (reverify) -- (confirmed);
\draw[separate] (confirmed) -- (refinegroup);

\node[note, draw=orange!65!black, rounded corners=1pt, fill=white] at (0,-6.85) {A refined candidate is not mixed into the original option group};

\end{tikzpicture}
}
\caption{Test-time role learning from homogeneous verified groups and the separate path for bounded refinement.}
\label{figtesttimelearning}
\end{figure}
